%
% Rejoinder for PA2534 resubmission
%
\newcommand{\theVersion}{1.0 -- October 2, 2026}
\documentclass[12pt,a4paper,twoside]{article}
\usepackage{times}
\usepackage{multirow}
\usepackage{hyperref}
\usepackage[T1]{fontenc}
\usepackage[utf8]{inputenc}
\usepackage[top=2.5cm, bottom=2.5cm, left=2.5cm, right=2.5cm]{geometry}
\usepackage[sort&compress]{natbib}
\setcitestyle{numbers,square,comma}
\usepackage{enumitem}
\setlist{topsep=1ex,itemsep=0.5ex,parsep=0pt,partopsep=0pt}
\setlength{\bibsep}{4pt}
\usepackage{longtable}
\usepackage{booktabs}
\usepackage{parskip}

\newcommand{\theCourse}{PA2534: Master's Thesis in Software Engineering}

\title{Rejoinder for degree project\\\vspace{1mm}\small{Version \theVersion}}
\author{\textsc{\theCourse}}
\date{October 2, 2026}

\begin{document}
\maketitle
\vspace*{-5mm}

\noindent
\begin{tabular}{|l|p{12.85cm}|}
\hline
Title         & Human-AI Collaboration for Decision-Making in Agile Sprint Planning: A Practitioner-Informed Conceptual Model \\
\hline
Author(s)     & Ungarala Sai Krishna Yashwanth; Mihir Singh Thakur \\
\hline
e-Mail(s)     & saun24@student.bth.se; mith24@student.bth.se \\
\hline
Supervisor(s) & Henry Edison \\
\hline\hline
Reviewer 1    & Michael Unterkalmsteiner (Examiner) \\
\hline
Reviewer 2    & (Anonymous second reviewer) \\
\hline
\end{tabular}

\section{Introduction}

This rejoinder responds to the examiner feedback received on the original submission.
The examiner identified the rubric items \textit{Method Selection and Application} and
\textit{Traceability and Validity} as requiring substantial revision. The examiner's
summary identified two required fixes:
(1) reporting the
research method application accurately and completely, and (2) accurately replicating
cited statements without paraphrasing that changes meaning.

The revised thesis responds to the specific comments raised by both reviewers. The
major changes are:

\begin{enumerate}
  \item A dedicated Background section (2.1) introducing sprint planning, the Scrum
        Master role, planning poker, anchoring bias, and algorithm aversion.
  \item Integration of the six missing papers identified by Reviewer 1.
  \item A rewritten research gap as three numbered specific gaps.
  \item An Alternative Research Methods section (3.2) systematically ruling out
        four alternatives.
  \item Phase-by-phase description of the data analysis (Phases 1--6 explicitly
        named and described).
  \item Complete quotation traceability: every quote now carries a label specifying
        the interview part, question number, and transcript timestamp (e.g., [Part 2,
        Vignette A, Q5, 9:54]).
  \item Silently cleaned quotations converted to explicitly labelled paraphrases.
  \item A second structured validation round via email (August--September 2026) with
        full verbatim responses in Appendix C.
  \item Correction of the bibliography errors identified in the feedback (paper mill
        output removed, arXiv status disclosed, and seven author/year corrections).
  \item Overclaiming language identified during revision was replaced with more
        descriptive language.
\end{enumerate}

\section{Response to reviewer reports}

\small{
\begin{longtable}{p{1cm}p{6.7cm}p{7cm}}
\toprule
\textbf{ID} & \textbf{Comment} & \textbf{Answer} \\
\midrule
\endhead
\bottomrule
\multicolumn{3}{l}{\footnotesize{R1 = Reviewer 1 (Examiner); R2 = Reviewer 2}}
\endlastfoot

R1.1 &
Chapters 1--2 are very short (not longer than the project plan). No introduction of
relevant concepts such as agile planning practices, planning poker, facilitator roles,
or team-level vs.\ individual decision making. &
A dedicated Background section (Section~2.1) has been added. It introduces five
concepts with empirical citations: sprint planning in Scrum (Mahnic 2011; \v{Z}\'a\v{c}ek
et al.\ 2025), the Scrum Master as facilitator (Shastri et al.\ 2021), planning poker
(Molokken-Ostvold \& Haugen 2007), anchoring bias (Kahneman 2011; Mohanani et al.\
2020), and algorithm aversion (Dietvorst et al.\ 2015). Chapter~2 now contains
approximately 2,570 words excluding references. \\

R1.2 &
Vague statements about cited sources; no critical assessment of methods and findings;
no synthesis; enumeration of papers. &
Every paper in Sections~2.2--2.5 is now described using a four-part pattern: research
approach, main finding, relevance to this thesis, remaining gap. \\

R1.3 &
Muneer [22] did not compare LLM-based to traditional approaches and does not conclude
that human oversight is necessary. &
Corrected in Section~2.3. Muneer is now described as an examination of different LLM
architectures for sprint planning support, without attributing conclusions the paper
does not make. \\

R1.4 &
Six relevant papers not cited (Cabrero-Daniel 2024, Tsihrintzis 2025, Martino 2025,
Alhubaishy 2026, Rahmahsari \& Sekti 2026, Brandas 2026). &
All six are now integrated in Sections~2.2--2.3 with the four-part description. \\

R1.5 &
HACO already cites Lee \& See in its 2020 trust section; the claimed synthesis ``While
HACO identifies trust\ldots Lee and See supply that theoretical foundation'' is false. &
Corrected in Section~2.5: ``The HACO trust dimension already draws on Lee and See's
foundational theory of calibrated trust.'' \\

R1.6 &
Research gap formulated as a cheap combination of criteria that only this thesis
satisfies. &
Rewritten as three numbered specific gaps in Section~2.6, each citing the specific
studies that leave the gap. \\

R1.7 &
``Perceived decision quality'' in RQ1a never defined or operationalised; no actual
planning decision observed. &
Removed from RQ1a throughout the thesis. The revised RQ1a asks only how practitioners
describe anchoring and algorithm aversion risk as affecting reasoning. \\

R1.8 &
Four aims / four research questions are essentially duplicates. &
A paragraph after the objectives (Section~1.2) distinguishes the two: objectives
describe research activities; RQs specify the analytical questions. \\

R1.9 &
Sprint planning does not exist in every Agile process. &
Corrected in Section~1.1 and the Background: the thesis now specifies sprint-based
Agile development (Scrum) and notes that Kanban and other methods do not use sprints. \\

R1.10 &
Method chapter has nice optics but is not a transparent description. Machine
transcription not disclosed; transcript accuracy not addressed. &
Section~3.4 now discloses: ``Transcription was completed using automated
speech-to-text software and reviewed by one researcher for accuracy prior to coding;
occasional recognition errors may remain.'' \\

R1.11 &
Protocol deviations (P5 only received Vignette A; wording adapted for P3) not
disclosed. &
Section~3.4 now includes a dedicated Protocol Deviations paragraph disclosing both
deviations. P5 is marked N/A for B1--B4 in Table~3.1; all Vignette B findings use
eight participants rather than nine. The adaptation for P3 was an ad hoc interviewer
decision; P2, who also had no direct AI tool experience, did not receive the same
adaptation. This inconsistency is acknowledged as a limitation. \\

R1.12 &
Saturation table (Table~3.1) makes no sense given the actual interview order (P5, P2,
P9, P7, P6, P8, P1, P4, P3); looks different from the May~10 draft. &
Table~3.1 renamed ``Code coverage across participants''; ``New codes introduced'' row
removed. Caption explains that P1--P9 are ordered by seniority not interview date.
Surrounding text states explicitly that saturation was not established: all interviews
were completed before coding began; recruitment stopped because of scheduling
constraints. The actual chronological interview order is stated in the text. \\

R1.13 &
``We continued recruitment until we reached thematic saturation [Braun \& Clarke]''
--- Braun \& Clarke do not discuss saturation. &
Citation removed. Sample size sentence now states ``a pragmatic recruitment goal
given the project scope and access constraints.'' \\

R1.14 &
Three Creswell page references (pp.~47--48, pp.~86--87, pp.~146--149) point to
unrelated content. &
All three removed. \\

R1.15 &
``Open coding with no predetermined categories'' contradicts the pre-existing
operational definitions in the March project plan. &
Section~3.5 now describes ``inductive thematic coding informed by sensitising
concepts'' in Phase~A, explicitly separated from deductive HACO mapping in Phase~B.
Interpretive criteria (anchoring, algorithm aversion, calibrated reliance) are
described as sensitising concepts. The actual codebook codes are A1--C5 (Appendix~B).
The phrase ``operational definitions used throughout coding'' has been replaced with
``interpretive criteria used to connect final codes to research questions.'' \\

R1.16 &
``Dismissal-without-engagement pattern associated with algorithm aversion risk'' is
one hell of a code --- apparently string-replaced everywhere including in the codebook,
which is part of the research data. &
Codebook definitions in Appendix~B now use short operational language only. The long
phrase is used only in the analysis text where codes are interpreted. \\

R1.17 &
Phase B is descriptive; does not show how HACO mapping challenged or improved
interpretation beyond organising codes into broad buckets. &
Sections~3.5 and~3.6 now explain two analytical insights from the HACO mapping:
(1) A3 and C4 revealed complementary structural safeguards; (2) the Learning Paradigm
dimension produced no codes, itself an analytical finding. \\

R1.18 &
Inter-coder agreement never stated. &
Section~3.5 Phase~4 now states: ``Formal inter-coder reliability was not calculated;
disagreements were resolved through discussion and codebook definitions were refined
until both researchers agreed.'' \\

R1.19 &
All quotations completely decontextualized: impossible to tell if Part~1, Part~2, or
Part~3. Quotations used as if describing real-world behaviour when at least some are
clearly Vignette A responses. &
Every quotation in Chapter~4 now carries a label specifying the interview part,
question number, and transcript timestamp, e.g., [Part~2, Vignette~A, Q5, 9:54].
A note at the start of Chapter~4 states that vignette responses represent hypothetical
reasoning, not observed behaviour. \\

R1.20 &
Case 2: P2's verbatim quotation is heavily edited, changing semantics from a hedged
comparison to a definitive statement used as evidence for trust in AI transparency. &
The P2 C1 quotation has been paraphrased with the disclosure ``(Paraphrase used; P2
qualified this as hypothetical because they had no direct AI-tool experience).'' See
Section~4.4.1. \\

R1.21 &
Validation evidence almost certainly heavily edited; may not be authentic; contains
information not given to participants (``2 out of 9 is low evidence''). Timeline is
impossible: interviews ended May~3; full model communicated and responses received
before May~10. &
A second structured validation round was conducted via email in August--September 2026
after receiving the examiner's feedback. Full verbatim responses from P6, P8 and P4
are in Appendix~C. The first round is now described as ``informal preliminary feedback
with limited traceability'' and is not treated as equivalent in evidential weight to
the documented second round. According to the authors' reconstruction, two refinements
were associated with the first round: adding sprint inertia and significant team
composition change as moderate risk conditions. No supervisor confirmation is claimed. \\

R1.22 &
The thesis treats identification of anchoring risk as evidence that anchoring bias
manifested. Awareness is not the same as observing the bias. &
Language throughout Chapter~4 revised. ``Bias manifested'' does not appear. The thesis
now says practitioners ``described'' or ``responses consistent with.'' \\

R1.23 &
The claim that the model ``emerged'' from open inductive coding is too strong: elements
like human final authority and transparency are part of HACO. &
Section~3.6 now states the model ``was constructed after analysis by combining
recurring participant accounts, the deductive HACO mapping, and the theoretical
concepts established before data collection.'' The word ``emerged'' has been replaced
throughout. \\

R1.24 &
Ref~[5] is likely a paper mill output. &
barua2025ai removed from the bibliography and all citations. \\

R1.25 &
Ref~[2] missing one author; Ref~[18] missing venue (arXiv preprint); TSE volumes 45
and 46 published in 2019 and 2020 not 2018. &
Seven bibliography corrections made (see Further Changes section for full list). \\

R1.26 &
Abstract labels results section ``Conclusions.'' &
Changed to ``Model.'' \\

R1.27 &
Page~2 formatting issues: vertical space in declaration, text running over margin,
bold-face ``Contact Information'' and ``University advisor'' in email line. &
Page~2 layout corrected: bold removed from ``Contact Information''; advisor line
repositioned; bottom table reformatted to fit within margins. \\

R1.28 &
Weird wording: ``adjust from an AI-suggested figure'' (abstract); ``We called those
spaces zones because that is what they are'' (Section~3.5); ``B2 fires explicitly in
P7 and P8'' (Section~4.5.2). &
All three corrected. Abstract now says ``treat an AI-suggested figure as a reference
point and adjust insufficiently from it.'' Section~3.5 no longer contains the ``We
called\ldots'' formulations. ``Fires explicitly'' replaced with ``was coded
explicitly.'' \\

R2.1 &
The report did not clearly describe the steps to develop the vignette scenarios as the
instrument to collect data. &
Section~3.4 now includes a paragraph explaining vignette development: grounded in
target bias mechanisms from the literature, calibrated following Aguinis \& Bradley's
ecological validity guidance, reviewed by both researchers against operational
definitions. \\

R2.2 &
The data analysis process and model construction process were not justified and
motivated. &
Phase-by-phase description of data analysis added (Phases~1--6 explicitly named).
Section~3.6 explains the three-phase model construction (Phase~A: inductive coding;
Phase~B: HACO mapping; Phase~C: model construction). \\

R2.3 &
``The interaction principles emerged because participants described practices.'' How
were the principles extracted from the practices? &
Section~3.6 now explains that a described practice was elevated to an interaction
principle when three conditions were met: multiple participants converged on the same
mitigation strategy; it was framed as a deliberate and repeatable behaviour; and it
addressed a specific bias risk identified in RQ1. \\

R2.4 &
Evidences are silent on how anchoring bias and aversion risk affect practitioner
reasoning and perceived quality. The quotes may not be relevant to support the claims. &
``Perceived decision quality'' removed from RQ1a. Every quote now tagged with
interview part, question, and timestamp. Quote tagging makes it possible to assess
whether each quote supports the claim it is used to support. Summary tables in
Chapter~4 show participant coverage and evidence type. Table~4.1 additionally reports
the applicable denominator for each finding. \\

\end{longtable}
}


\section{Response to further feedback}

The following comments were raised by the supervisor (Henry Edison) when reviewing
the revised draft before resubmission.

\small{
\begin{longtable}{p{1cm}p{6.7cm}p{7cm}}
\toprule
\textbf{ID} & \textbf{Comment} & \textbf{Answer} \\
\midrule
\endhead
\bottomrule
\multicolumn{3}{l}{\footnotesize{Sup = Supervisor}}
\endlastfoot

Sup1 &
Chapters~1 and 2 still too short; no introduction of relevant concepts. &
Addressed --- see R1.1 above. \\

Sup2 &
Why did you not follow the vignette methodology directly given that you admit it has
its own methodology? &
The study used vignettes as an elicitation technique embedded within semi-structured
qualitative interviews, not as a standalone experimental vignette study. Section~3.1
now makes this boundary explicit: Part~1 elicited accounts of real practice; Part~2
elicited reasoning about hypothetical situations; Part~3 elicited trust conditions
and preferred future arrangements. We followed Aguinis \& Bradley's ecological
validity guidance for scenario design but did not manipulate experimental conditions
or make causal comparisons between groups. \\

Sup3 &
P2 has no AI tool experience (like P3); did you adapt the wording for P2? What
adaptation did you actually make for P3? &
P2 received the standard first-person vignette wording. No adaptation was made for
P2. P3 received an ad hoc third-person introduction: ``Since your team has not used
an AI planning tool directly, I will frame these scenarios slightly differently.
Imagine your team is using a planning dashboard or some automated tool\ldots'' The
numerical content remained unchanged. The inconsistency between P2 and P3 is now
acknowledged explicitly in Section~3.4 as a limitation. \\

Sup4 &
How does the absence of Vignette B from P5 affect the results? &
P5 is now marked N/A for B1--B4 in Table~3.1 with the note ``not elicited because
Vignette B was not administered.'' All Vignette B findings use an analysable
denominator of eight participants rather than nine. This is acknowledged as a
limitation. \\

Sup5 &
There are different sets of concepts in Chapter~3 and Appendix~B. &
Three levels of concepts are now distinguished in Section~3.5: (1) a priori
sensitising concepts (anchoring, algorithm aversion, calibrated reliance); (2) the
actual transcript-level codes A1--C5 in Appendix~B; (3) HACO dimensions used
deductively in Phase~B. The phrase ``operational definitions used throughout coding''
has been replaced with ``interpretive criteria used to connect final codes to research
questions.'' \\

Sup6 &
The quotes in Table~3.3 do not clearly show the connection to the concepts in
Chapter~3; longer quotes with more context would help. &
Table~3.3 has been revised: it is described as a representative evidence-chain table
rather than the complete coding evidence. Where automated transcripts were unclear,
entries have been converted to labelled paraphrases. The B1 row has been replaced with
P4's real incident (23:30--24:24), which provides clearer analytical grounding than
the previous P2 entry. \\

Sup7 &
I still do not understand how you came up with the model components in Table~3.4.
What is the purpose of the model? &
Section~3.6 now explains the three-phase construction: Phase~A codes, Phase~B HACO
mapping, Phase~C synthesis. Task characteristics informed zone definitions and
number-withholding principle; Teaming characteristics informed facilitator mediation,
examine-both-signals, and override justification; Trust informed transparency and
risk conditions. The purpose is stated: a practitioner-oriented structure for deciding
what AI may do, how AI information reaches the planning discussion, which decisions
must remain human, which safeguards govern disagreement, and which conditions indicate
elevated risk. \\

Sup8 &
``Their responses were discussed with the thesis supervisor in a Teams meeting prior
to the final submission.'' When did we discuss this? What was the output? Where are
the full first-round responses? &
This sentence has been removed. The first round is now described as ``informal
preliminary feedback with limited traceability; a complete written record is not
available.'' According to the authors' reconstruction, it was associated with two
refinements: adding sprint inertia and significant team composition change as moderate
risk conditions. These reconstructed outputs are reported with an explicit traceability
caveat, and no supervisor confirmation is claimed. The second round (email,
August--September 2026) produced the verifiable record in Appendix~C. \\

Sup9 &
What was the purpose of the second validation? What new things did it bring? &
The purpose was to produce a credible written record that could be examined directly.
Concrete outputs: two new checklist items added (item~12: unacknowledged trust
erosion from P6; item~13: short-term availability blind spots from P8); item~9
wording changed (``developer override asserted without articulated basis''); P6
extended the human authority principle; P4's response revealed that some principles
are aspirational, leading to framing the model as a proposed standard rather than a
description of current practice. \\

Sup10 &
RQ3 has a team angle but Section~3.7.3 leaves it as future work. How can findings
from individuals answer the team angle? &
RQ3 has been narrowed to: ``What practitioner-informed principles can be proposed for
structuring Human-AI collaboration in Agile sprint planning while preserving human
oversight?'' Section~3.7.3 explicitly states that team-level effectiveness, collective
acceptance and behavioural outcomes remain untested and require future team-level
observation or intervention studies. \\

Sup11 &
How did you decide on the ``Strength'' labels in Table~4.1? &
The Strength column has been replaced by a Denominator column. The table now reports
participant count, applicable denominator (n=8 for Vignette~B findings, n=9 for
others), and evidence type. \\

Sup12 &
Check consistency of all reported numbers across chapters. &
A systematic audit has been performed. All Vignette~B findings now use n=8.
Principle~4 evidence states 2/8. The checklist consistently shows 13 conditions.
P5 is N/A for B1--B4 throughout. \\

Sup13 &
Explain how Figure~4.1 can be used. &
Figure~4.1 caption and surrounding paragraph revised. The figure now functions as a
process map: before planning, the facilitator uses Zone~1 to review AI analysis
privately; during planning, Zone~2 shows how signals can be introduced without
exposing the team to an anchor; Zone~3 identifies human-only decisions. The five
principles govern information flow between zones; the risk checklist indicates when
additional safeguards are required. \\

\end{longtable}
}


\section{Further changes}

The following changes were made proactively to improve the thesis, not in direct
response to a specific reviewer comment.

\begin{enumerate}

  \item \textbf{Research Scope section added} (Section~1.6). Explicitly states what
        the thesis does not examine: non-sprint-based Agile methods, AI tool accuracy,
        team-level outcomes, other Agile ceremonies.

  \item \textbf{Alternative Research Methods section added} (Section~3.2). Four
        alternatives systematically ruled out: controlled experiment, quasi-experiment,
        direct observation, artefact-based case study.

  \item \textbf{Implications for Research added} (Section~5.7). Four specific
        directions for future empirical work.

  \item \textbf{Implications for Practice added} (Section~5.8). Three
        practitioner-oriented recommendations.

  \item \textbf{Three summary tables added to Results} (after RQ1a, RQ1b, RQ2).
        Table~4.1 reports participant count, applicable denominator, and evidence type.
        The RQ1b and RQ2 summary tables report participant coverage and evidence type.

  \item \textbf{Two additional evidence corrections were made during the final audit}.
        First, P2 was recoded from B1 to B4. Table~3.3 now assigns B1 to P4's recounted
        real incident (23:30--24:24), and the Results state that no participant directly
        demonstrated B1 in their own Vignette~B response. Second, a composite P7
        quotation that combined two separate incidents was removed; P7's evidence now
        appears only as correctly attributed quotations and paraphrases.

  \item \textbf{Two new Discussion sub-sections added}: Section~5.3 (Convergence of
        Evidence Across Interview Parts) distinguishes stronger from weaker evidential
        support; Section~5.4 replaces the prior experience-moderation claim with an
        honest statement that the study did not establish a directional relationship
        between experience and bias calibration.

  \item \textbf{Bibliography corrections}: seven entries corrected for author names,
        publication years, or missing venue:
        \begin{itemize}[noitemsep]
          \item aman2025systematic: four authors corrected; pages 189--217; DOI added.
          \item klieger2024chatcollab: three author names corrected (Charis Charitsis,
                Miroslav Suzara, Sierra Wang); arXiv preprint disclosed.
          \item emmanouilidis2025co: three author names corrected (Jessica Zotelli,
                Katharina Hengel, Jos Bokhorst); DOI added.
          \item hamza2024human: three author names corrected (Muhammad Hamza, Muhammad
                Azeem Akbar, Tahsinur Rahman); DOI added.
          \item \v{Z}\'{a}\v{c}ek 2025: two first names corrected (Miroslav
                \v{Z}\'{a}\v{c}ek, Ad\'{e}la Hamplov\'{a}).
          \item mohanani2018cognitive: year corrected to 2020 (TSE vol.~46).
          \item choetkiertikul2018deep: year corrected to 2019 (TSE vol.~45).
        \end{itemize}

  \item \textbf{GenAI disclosure updated}: Claude (Anthropic) named specifically as
        the tool used for language assistance and vignette scenario development.

  \item \textbf{Consent priming acknowledged as limitation} (Section~5.5): the
        consent information sheet disclosed the study's focus on cognitive bias before
        interviews, which may have primed participant responses.

\end{enumerate}

\end{document}